% TEMPLATE-MILC-v3.tex - plantilla maestra UNICA de las guias MILC v3.
%
% La usan las tres familias de guias, cada una con su driver:
%   · Grados 6-11  -> scripts/build-guias-g11.py
%   · Bebras       -> scripts/build-guias-bebras.py
%   · Semillero    -> scripts/build-guias-semillero.py
%
% Antes habia tres copias casi identicas (~400 lineas iguales, 6 distintas):
% cada mejora habia que aplicarla tres veces, y en la practica no se hacia
% (el motor de recursos vivia solo en dos de ellas). Lo que varia entre
% familias esta parametrizado en 6 placeholders que cada driver llena
% (se nombran aqui SIN su sintaxis real: el builder reemplaza por texto plano,
% tambien dentro de los comentarios, y un valor multilinea rompe el preambulo):
%
%   HEADER_IZQ        encabezado izquierdo de pagina
%   HEADER_DER        encabezado derecho de pagina
%   PORTADA_ETIQUETA  rotulo sobre el numero grande (GRADO/MOMENTO/MODULO)
%   PORTADA_SERIE     linea de serie de la portada
%   PORTADA_ID        identificador de la guia en la portada
%   PORTADA_CIRCULO   el circulito de grado (°), o vacio si no aplica
%
% El resto de claves las documenta content/guias/_SCHEMA.md. El script valida
% que no queden placeholders sin reemplazar antes de escribir el archivo final.

\documentclass[11pt,letterpaper]{article}
\usepackage[margin=2.0cm]{geometry}
\usepackage[table]{xcolor}
\usepackage{fontspec}
\usepackage[spanish]{babel}
\usepackage{tikz}
% Librerías TikZ para los diagramas de las guías (bloque `recursos:`):
%  · babel  → babel en español vuelve activos caracteres como «>», y TikZ los usa
%             en sus opciones (>=stealth, ->). Sin esta librería, cualquier
%             diagrama con flechas revienta con «\language@active@arg> has an
%             extra }». Debe ir ANTES de que se dibuje nada.
%  · positioning        → sintaxis «below left=2pt and 3pt».
%  · decorations.pathreplacing → llaves ({brace}) para señalar rangos.
\usetikzlibrary{babel,positioning,decorations.pathreplacing}
\usepackage{graphicx}
% Circuitos: el trabajo de electrónica se hace hoy con micro:bit y sus sensores
% integrados (MakeCode, sin cableado), así que no se necesitan esquemáticos.
% Si algún día entran montajes con protoboard, basta descomentar esta línea:
% los diagramas .tex/.tikz declarados en `recursos:` ya se incrustan solos.
% \usepackage{circuitikz}
\usepackage[most]{tcolorbox}
\usepackage{tabularx}
\usepackage{array}
\usepackage{fancyhdr}
\usepackage{titlesec}
\usepackage[document]{ragged2e}  % bandera a la derecha en todo el cuerpo
\usepackage{enumitem}
\usepackage{microtype}

% Helvetica Neue no trae la flecha U+2192, asi que cada «→» escrito literalmente
% en el YAML se compilaba a NADA, en silencio (462 flechas en 61 guias: rutas del
% tipo «paleta → tipografia → lockup» salian rotas). Mapeamos el caracter a su
% version matematica, que si tiene glifo. Asi el autor puede seguir escribiendo
% «→» comodamente en el YAML.
\usepackage{newunicodechar}
\newunicodechar{→}{\ensuremath{\rightarrow}}
% Mismo problema con los simbolos de marca y vinetas: el «✓» de «una palomita ✓»
% salia como cuadrito vacio en el PDF de todas las guias que lo usaban.
\usepackage{pifont}
\newunicodechar{✓}{\ding{51}}
\newunicodechar{✗}{\ding{55}}
\newunicodechar{✔}{\ding{52}}
\newunicodechar{•}{\textbullet}
\newunicodechar{≥}{\ensuremath{\geq}}
\newunicodechar{≤}{\ensuremath{\leq}}

\setmainfont{Helvetica Neue}
\setsansfont{Helvetica Neue}
\setlength{\parindent}{0pt}
\setlength{\parskip}{6pt}
% Sin particion de palabras: en 6.º-7.º una palabra cortada por guion al final
% de linea («Comuni-cacion») frena la lectura mucho mas que una linea corta.
\hyphenpenalty=10000
\exhyphenpenalty=10000
\pretolerance=2000
\tolerance=3000
\setlength{\emergencystretch}{3em}
\linespread{1.10}
\renewcommand{\arraystretch}{1.25}
\setlist{leftmargin=5mm,itemsep=1mm,topsep=1mm}
\newcolumntype{Y}{>{\RaggedRight\arraybackslash}X}

\definecolor{milcMagenta}{HTML}{E6007E}
\definecolor{milcVino}{HTML}{5A0038}
\definecolor{milcTurquesa}{HTML}{0093A5}
\definecolor{milcVerde}{HTML}{58A923}
\definecolor{milcMostaza}{HTML}{E5B400}
\definecolor{milcOcre}{HTML}{B86600}
\definecolor{milcNegro}{HTML}{241816}
\definecolor{milcGris}{HTML}{F7F7F4}
\definecolor{milcLinea}{HTML}{E8E2DD}
\definecolor{milcGradePrimary}{HTML}{7C3AED}
\definecolor{milcGradeDark}{HTML}{5B21B6}
\definecolor{milcGradeSoft}{HTML}{EDE9FE}
\definecolor{milcGradeAccent}{HTML}{CFE7FF}
\definecolor{milcGradeText}{HTML}{FFFFFF}
\color{milcNegro}

\pagestyle{fancy}
\fancyhf{}
\lhead{\small Semillero de Investigación · Astronomía y ciencia ciudadana}
\rhead{\small Módulo 3 · Praxis · MILC v3}
\cfoot{\small\thepage}
\renewcommand{\headrulewidth}{0pt}

\newtcolorbox{titlebox}[1]{
  enhanced,colback=#1,colframe=#1,arc=7mm,boxrule=0pt,
  left=7mm,right=7mm,top=3mm,bottom=3mm,
  before skip=8pt,after skip=8pt,
  fontupper=\sffamily\bfseries\Large\color{white}
}

\newtcolorbox{softbox}[3]{
  enhanced,colback=#2,colframe=#1,borderline west={4pt}{0pt}{#1},
  arc=2mm,boxrule=.4pt,left=4mm,right=4mm,top=3mm,bottom=3mm,
  before skip=6pt,after skip=8pt,title={#3},coltitle=white,
  colbacktitle=#1,fonttitle=\sffamily\bfseries,fontupper=\RaggedRight,
  attach boxed title to top left={xshift=3mm,yshift=-2mm},
  boxed title style={arc=2mm, boxrule=0pt}
}

\newtcolorbox{drawbox}[2]{
  enhanced,colback=white,colframe=#1,arc=4mm,boxrule=1pt,
  left=5mm,right=5mm,top=4mm,bottom=4mm,before skip=8pt,after skip=8pt,
  title={#2},coltitle=white,colbacktitle=#1,fonttitle=\sffamily\bfseries,
  fontupper=\RaggedRight,
  attach boxed title to top left={xshift=4mm,yshift=-2mm},
  boxed title style={arc=3mm, boxrule=0pt}
}

\newtcolorbox{infoband}[1]{
  enhanced,colback=#1!8,colframe=#1!50,arc=2mm,boxrule=.5pt,
  left=4mm,right=4mm,top=2mm,bottom=2mm,before skip=4pt,after skip=4pt,
  fontupper=\small\RaggedRight
}

% Puente narrativo entre secciones (contrato editorial Conéctate).
% Da continuidad: "Acabas de X. Ahora vas a Y porque Z."
% Visualmente: banda izquierda en color del grado, texto en itálica pequeña.
\newtcolorbox{puentebox}{
  enhanced,colback=milcGradeSoft,colframe=milcGradeDark,
  borderline west={3pt}{0pt}{milcGradeDark},
  arc=0mm,boxrule=0pt,left=5mm,right=4mm,top=2.5mm,bottom=2.5mm,
  before skip=4pt,after skip=8pt,
  fontupper=\itshape\small\color{milcNegro}\RaggedRight
}
% La flecha va en modo matematico a proposito: Helvetica Neue NO tiene el glifo
% U+2192, asi que \textrightarrow se compilaba a NADA («Missing character: There
% is no → in font Helvetica Neue») y el puente salia sin flecha. En modo
% matematico la flecha viene de la fuente matematica y si se dibuja.
\newcommand{\puente}[1]{%
  \begin{puentebox}%
    \textcolor{milcGradeDark}{$\rightarrow$}\hspace{2mm}#1%
  \end{puentebox}%
}

\newcommand{\linead}[1]{\noindent\color{milcLinea}\rule{#1}{0.5pt}\color{milcNegro}}
\newcommand{\respuesta}[1]{\par\vspace{2mm}\foreach \n in {1,...,#1}{\noindent\color{milcLinea}\rule{\linewidth}{0.5pt}\par\vspace{5mm}}\color{milcNegro}}
\newcommand{\checkbox}{\textcolor{milcGradeDark}{\fbox{\phantom{x}}}\hspace{5pt}}

\newcommand{\phasecard}[4]{
\begin{tcolorbox}[enhanced,colback=#3,colframe=#2,arc=3mm,boxrule=.5pt,height=3.05cm,valign=center,halign=center,left=1.5mm,right=1.5mm,top=2mm,bottom=2mm]
{\sffamily\bfseries\fontsize{22}{24}\selectfont\textcolor{#2}{#1}}\par
{\sffamily\bfseries\small #4}
\end{tcolorbox}
}

% ── Piezas del rediseno 2026 (legibilidad 6.º-11.º) ───────────────────
% Banda de estacion: reemplaza el titlebox macizo. Titulo a la izquierda,
% dato practico (tiempo/modalidad) a la derecha, en una sola linea.
\newcommand{\milcestacion}[2]{%
\begin{tcolorbox}[enhanced,colback=#1,colframe=#1,arc=2mm,boxrule=0pt,
  left=5mm,right=5mm,top=2.6mm,bottom=2.6mm,before skip=11pt,after skip=7pt]
{\sffamily\bfseries\fontsize{15}{18}\selectfont\color{white}#2}
\end{tcolorbox}}

% Tarjeta de paso numerada: sustituye las tablas «Pilar / Aplicacion», que
% eran formato de consulta docente, no de estudiante. El numero va en un
% circulo sobre el borde para que la secuencia se lea de un vistazo.
\newcommand{\milcpaso}[3]{%
\begin{tcolorbox}[enhanced,colback=white,colframe=milcLinea,arc=1.5mm,boxrule=.9pt,
  left=14mm,right=4mm,top=3mm,bottom=3mm,before skip=4pt,after skip=4pt,
  fontupper=\RaggedRight,
  overlay={\node[anchor=north west,circle,fill=#2,text=white,inner sep=0pt,
    minimum size=7.4mm,font=\sffamily\bfseries\small]
    at ([xshift=3.6mm,yshift=-3.2mm]frame.north west) {#1};}]
#3
\end{tcolorbox}}

% Casilla marcable de tamano util con lapiz (la anterior era \fbox{\phantom{x}},
% de alto variable segun el texto que la seguia).
\newcommand{\milcmarca}{\framebox[3.8mm]{\rule{0pt}{3.4mm}}}

% Fila marcable de lista suelta (checks de la estacion 1).
\newcommand{\milccheckrow}[1]{%
\par\vspace{1.8mm}\noindent
\makebox[6.4mm][l]{\raisebox{-.55mm}{\textcolor{milcNegro}{\milcmarca}}}%
\parbox[t]{\dimexpr\linewidth-6.4mm\relax}{\RaggedRight #1}\par}

% Celda marcable dentro de la rubrica (tabla de autoevaluacion). Misma
% sangria francesa, pero pensada para vivir dentro de una columna X.
\newcommand{\milcopcion}[1]{%
\makebox[5.2mm][l]{\raisebox{-.55mm}{\textcolor{milcNegro}{\milcmarca}}}%
\parbox[t]{\dimexpr\linewidth-5.2mm\relax}{\RaggedRight #1}}

% ── Recursos visuales de la guía ──────────────────────────────────────
% Declarados en el bloque `recursos:` del YAML y emitidos por el builder.
% \guiaFigura   → imagen rasterizada o PDF (foto, carta celeste, montaje).
% \guiaDiagrama → fuente TikZ/circuitikz (.tex) incrustada como vector:
%                 es la vía para circuitos y esquemáticos editables.
% #1 = ruta relativa al .tex · #2 = pie (puede ir vacío)
\newcommand{\guiaFigura}[2]{%
\begin{center}
\includegraphics[width=0.88\linewidth,height=0.38\textheight,keepaspectratio]{#1}\\[1.5mm]
{\footnotesize\itshape\textcolor{milcNegro!75}{#2}}
\end{center}
}

\newcommand{\guiaDiagrama}[2]{%
\begin{center}
\input{#1}\\[1.5mm]
{\footnotesize\itshape\textcolor{milcNegro!75}{#2}}
\end{center}
}

\begin{document}
\thispagestyle{empty}

% ── Portada ───────────────────────────────────────────────────────────
% Antes: un rectangulo de color a pagina completa. Se veia bien en pantalla,
% pero en la fotocopiadora del colegio se comia el toner y salia gris sucio.
% Ahora la banda ocupa el tercio superior y el resto va en blanco.
\begin{tikzpicture}[remember picture,overlay]
  \path[fill=milcGradePrimary]
    (current page.north west) rectangle ([yshift=-10.6cm]current page.north east);

  \node[anchor=north west,text=milcGradeText] at ([xshift=2.0cm,yshift=-1.55cm]current page.north west)
    {\sffamily\bfseries\fontsize{12}{14}\selectfont MÓDULO};
  \node[anchor=north west,text=milcGradeText,opacity=.85] at ([xshift=2.0cm,yshift=-2.35cm]current page.north west)
    {\sffamily\bfseries\fontsize{8.5}{10}\selectfont SEMILLERO DE INVESTIGACIÓN · CONECTATE};
  \node[anchor=north east,text=milcGradeText,opacity=.85,align=right] at ([xshift=-2.0cm,yshift=-1.55cm]current page.north east)
    {\sffamily\bfseries\fontsize{9}{12}\selectfont I.E. Sor María Juliana\\Cartago, Valle del Cauca};

  \node[anchor=west,text=milcGradeText] (gradonum) at ([xshift=1.85cm,yshift=-7.1cm]current page.north west)
    {\sffamily\bfseries\fontsize{112}{112}\selectfont 3};
  % El circulito de grado (°) solo aplica a las guias de grado («11°»); en
  % Bebras y Semillero el numero grande es un momento/modulo, no un grado.
  % Se posiciona relativo al nodo `gradonum`, no en coordenadas absolutas.
  

  \node[anchor=west,text=milcGradeText,text width=11.4cm,align=left] at ([xshift=1.5cm,yshift=0.5cm]gradonum.east)
    {
     {\sffamily\bfseries\fontsize{9}{11}\selectfont\MakeUppercase{Módulo 3 · Fase MILC: Praxis · 3 h de clase}}\\[3mm]
     {\sffamily\bfseries\fontsize{21}{25}\selectfont\setlength{\baselineskip}{27pt}Del ojo al instrumento ---\\[4pt]un fotómetro con micro:bit\\[4pt]para medir el cielo}\\[3.5mm]
     {\sffamily\bfseries\fontsize{15}{18}\selectfont Astronomía y ciencia ciudadana}
    };
\end{tikzpicture}

\vspace*{9.4cm}
\vfill

{\sffamily\bfseries\fontsize{8}{10}\selectfont\textcolor{milcGradeDark}{LO QUE TE LLEVAS HOY}}

\begin{tcolorbox}[enhanced,colback=white,colframe=milcNegro,arc=2mm,boxrule=1.6pt,
  left=5mm,right=5mm,top=4mm,bottom=4mm,before skip=3pt,after skip=12pt,
  fontupper=\RaggedRight\fontsize{11}{15.5}\selectfont]
Tu \textbf{fotómetro de contaminación lumínica con micro:bit}, construido y calibrado: el programa en MakeCode (con su \textbf{pseudocódigo} previo y su \textbf{tabla de umbrales} justificada), probado en \textbf{5 escenarios} (esperado / real). Más tu \textbf{tabla de contraste}: la lectura del instrumento frente a la estimación \emph{a ojo} de tu magnitud límite del módulo 1 ---¿coinciden?, ¿dónde no?, ¿por qué?
\end{tcolorbox}

\begin{tcolorbox}[enhanced,colback=milcGris,colframe=milcGris,arc=2mm,boxrule=0pt,
  left=5mm,right=5mm,top=3.5mm,bottom=3.5mm,after skip=12pt]
{\sffamily\bfseries\fontsize{8}{10}\selectfont\textcolor{milcNegro!55}{ESTA GUÍA ES DE}}\\[3mm]
\begin{tabularx}{\linewidth}{X p{3.2cm} p{3.2cm}}
\linead{\linewidth} & \linead{3.0cm} & \linead{3.0cm}\\[-1mm]
{\fontsize{8.5}{10}\selectfont\textcolor{milcNegro!55}{Nombre completo}} &
{\fontsize{8.5}{10}\selectfont\textcolor{milcNegro!55}{Grupo}} &
{\fontsize{8.5}{10}\selectfont\textcolor{milcNegro!55}{Fecha}}\\
\end{tabularx}
\end{tcolorbox}

\vfill

\noindent\textcolor{milcLinea}{\rule{\linewidth}{1pt}}\\[2mm]
{\sffamily\bfseries\fontsize{9.5}{12}\selectfont Autor: Dr. Álvaro Cárdenas Orozco}\\[1mm]
{\sffamily\fontsize{8.5}{11}\selectfont\textcolor{milcNegro!55}{Metodología MILC · apertura ancestral · pensamiento computacional · triángulo Dussel–estoicismo–Floridi}}

\newpage

\section*{Apertura: Del ojo al instrumento --- un fotómetro con micro:bit para medir el cielo}

\begin{softbox}{milcGradeDark}{milcGradeSoft}{Saber ancestral}
\textbf{En la galería de Cartago, durante siglos se vendió «al ojo».} La casera experimentada tomaba un puñado de fríjol, lo sopesaba en la mano y decía \emph{``esto es una libra''} ---y casi siempre acertaba. El buen ojo del oficio es real: años de práctica calibran la mano mejor de lo que uno cree. Pero \textbf{el ojo tiene dos límites}. Uno: \emph{cansa y varía} ---la misma mano pesa distinto en la mañana y al final del día. Otro, más serio: \emph{no se puede compartir ni comprobar}. Si el cliente decía \emph{``eso no es una libra''}, no había forma de resolver la discusión\ldots hasta que llegó \textbf{la pesa}. La balanza no vino a despreciar el ojo de la casera: vino a \textbf{volverlo exacto, justo y verificable}. Ahora los dos miran el mismo número, y nadie puede hacer trampa. Por algo las culturas antiguas guardaban las \textbf{pesas justas} como un asunto sagrado: un instrumento de medida bien calibrado es, antes que tecnología, un \textbf{acuerdo de justicia}. El ojo estima; el instrumento \textbf{da fe}.
\end{softbox}

\begin{softbox}{milcTurquesa}{milcGris}{Saber contemporáneo}
En el módulo 1 estimaste el brillo del cielo \emph{a ojo} (contando estrellas: tu magnitud límite). Hoy le pones una \textbf{pesa a la luz}. El \textbf{micro:bit} ---un computador del tamaño de una galleta--- trae un \textbf{sensor de luz integrado}: su matriz de 25 LED, además de mostrar, \textbf{mide} cuánta luz le llega. En MakeCode se lee con \texttt{input.lightLevel()}, que entrega un número de \textbf{0 (oscuridad) a 255 (mucha luz)}. \textbf{Sin cablear nada}: el instrumento ya está dentro. Pero un número crudo no dice nada por sí solo: hay que \textbf{calibrarlo}. Calibrar es \textbf{decidir umbrales} ---por debajo de tal valor «cielo oscuro», por encima «cielo con luz»--- midiendo condiciones conocidas, igual que la casera ajustaba su pesa con un peso patrón. Con eso construyes un \textbf{fotómetro de contaminación lumínica}: lo apuntas al cielo y te dice, con un número \emph{repetible y comparable}, cuánta luz de fondo hay. \emph{Ojo:} el micro:bit no compite con un observatorio; sirve para \textbf{comparar} lugares y noches (tu patio contra una esquina más oscura), que es justo lo que investiga la ciencia ciudadana de cielos.
\end{softbox}

\begin{softbox}{milcMostaza}{milcGris}{Pregunta-puente}
\textit{La casera sabía sopesar una libra, pero la pesa hizo su juicio \textbf{exacto y comprobable}; tú sabes estimar un cielo oscuro, pero\ldots \textbf{¿qué gana tu observación cuando el «me pareció oscuro» se vuelve un «medí 42» que otro puede repetir?}}
\end{softbox}

\begin{softbox}{milcVerde}{milcGris}{Saber y saber hacer}
Al terminar la sesión: (1) podrás \textbf{identificar} dónde el ojo estima bien y dónde hace falta un instrumento que dé fe; (2) sabrás \textbf{explicar} cómo el micro:bit mide luz y qué es calibrar con umbrales; (3) podrás \textbf{crear} y calibrar un fotómetro en MakeCode, probándolo en 5 escenarios; (4) habrás \textbf{evaluado} tu instrumento contrastando su lectura con tu estimación a ojo del módulo 1.
\end{softbox}



\small
\begin{tabularx}{\textwidth}{>{\bfseries}p{3.0cm}Y}
\rowcolor{milcGradePrimary!12}Autor & Dr. Álvaro Cárdenas Orozco\\
DBA & Formula preguntas investigables, produce evidencia con método y comunica hallazgos reconociendo sus límites (investigación estudiantil STEM, transversal 6°--11°)\\
Referentes & Orlove, Chiang \& Cane (2000, Nature) · RECA · NASA-IASC · Saberes campesinos y Quimbaya del Norte del Valle\\
Producto final & Tu \textbf{fotómetro de contaminación lumínica con micro:bit}, construido y calibrado: el programa en MakeCode (con su \textbf{pseudocódigo} previo y su \textbf{tabla de umbrales} justificada), probado en \textbf{5 escenarios} (esperado / real). Más tu \textbf{tabla de contraste}: la lectura del instrumento frente a la estimación \emph{a ojo} de tu magnitud límite del módulo 1 ---¿coinciden?, ¿dónde no?, ¿por qué?\\
\end{tabularx}
\normalsize

\puente{En el módulo 1 miraste con el ojo; en el 2 ordenaste los datos. Hoy \textbf{construyes el instrumento}: un fotómetro con micro:bit que convierte tu «me pareció oscuro» en un número que otro puede comprobar.}

\milcestacion{milcGradeDark}{Ruta MILC: así vamos a aprender}
\begin{tabularx}{\textwidth}{>{\centering\arraybackslash}X>{\centering\arraybackslash}X>{\centering\arraybackslash}X>{\centering\arraybackslash}X}
\phasecard{1}{milcVerde}{milcGris}{Escucha\\[-1mm]\small Observo, escucho y pregunto.} &
\phasecard{2}{milcTurquesa}{milcGris}{Entender\\[-1mm]\small Sistematizo y ordeno.} &
\phasecard{3}{milcMagenta}{milcGris}{Hacer\\[-1mm]\small Construyo y aplico.} &
\phasecard{4}{milcVino}{milcGris}{Cierre\\[-1mm]\small Evalúo y reflexiono.}
\end{tabularx}


\puente{Antes de construir, reconoce cuándo el ojo basta y cuándo engaña. Tu propia magnitud límite del módulo 1 es el punto de partida: ¿qué tan confiable fue?}

\milcestacion{milcVerde}{Estación 1 de 3 · Escucha}

\begin{softbox}{milcVerde}{milcGris}{Escena para escuchar}
\textbf{Actividad 1 · IDENTIFICA --- ¿Cuándo basta el ojo y cuándo engaña?} (35 min · individual + grupo). \textbf{Momento A (10 min) --- Prueba tu ojo.} El profe muestra 3 pares de cosas (dos grises casi iguales, dos sonidos, dos pesos en la mano) y cada quien estima cuál es «más». Se revela la respuesta con un instrumento (regla, balanza). \emph{¿Acertó tu ojo? ¿En qué falló?}. \textbf{Momento B (15 min) --- El caso de la pesa.} El profe presenta el paso de «vender al ojo» a «vender con pesa» y la pregunta: \emph{¿la pesa reemplazó el ojo de la casera, o lo hizo confiable para todos?}. Discutan: ¿qué gana una medida cuando \textbf{se puede compartir y comprobar}? \textbf{Momento C (10 min) --- Tu ojo del módulo 1.} Saca tu magnitud límite (las estrellas que contaste). \emph{¿La contarías igual otra noche? ¿Otro compañero contaría lo mismo?}. \textbf{En el cuaderno:} en qué acertó y en qué falló tu ojo + una frase sobre qué le da un instrumento a una observación (repetible · comparable · comprobable).
\end{softbox}

{\sffamily\bfseries\fontsize{8}{10}\selectfont\textcolor{milcNegro!55}{MARCA CUANDO LO HAYAS HECHO}}
\milccheckrow{Estimé a ojo y comparé con la medida real de un instrumento.}
\milccheckrow{Puedo decir qué gana una medida cuando se puede compartir y comprobar.}
\milccheckrow{Reconocí un límite de mi estimación a ojo del módulo 1.}

\begin{infoband}{milcVerde}


\textbf{Lo que va al cuaderno (Actividad 1 · IDENTIFICA).} Título: «Actividad 1 --- Cuándo basta el ojo». \textbf{Formato:} tabla de mis estimaciones a ojo vs la medida real + una frase sobre qué aporta un instrumento (repetible, comparable, comprobable). \textbf{Extensión:} 8 líneas. \textbf{Sabes que terminaste cuando} puedes nombrar \textbf{un caso} donde tu ojo se equivocó y un instrumento lo habría resuelto.
\end{infoband}

\puente{Ya sabes por qué hace falta un instrumento. Ahora entiende cómo mide luz el micro:bit y por qué un número crudo no sirve sin \textbf{umbrales calibrados}.}

\milcestacion{milcTurquesa}{Estación 2 de 3 · Entender}

\begin{softbox}{milcTurquesa}{milcGris}{Lo que vas a entender}
\textbf{Actividad 2 · EXPLICA --- Cómo mide luz el micro:bit y qué es calibrar} (55 min · grupo + individual). Ya sabes por qué hace falta un instrumento. Ahora entiende \textbf{cómo} el micro:bit mide luz, por qué un número crudo no basta, y cómo se \textbf{diseña la lógica} antes de tocar un bloque.
\end{softbox}

\milcpaso{1}{milcTurquesa}{\textbf{Pilar 1 --- El instrumento ya está dentro.} El micro:bit trae un \textbf{sensor de luz integrado}: su matriz de 25 LED, que normalmente muestra íconos, también \textbf{mide} la luz que le llega (los LED, al revés, funcionan como pequeños detectores). En MakeCode se lee con \texttt{input.lightLevel()} y devuelve un número de \textbf{0 a 255}: 0 en oscuridad total, 255 con mucha luz encima. \textbf{No hay que cablear ni comprar nada:} el fotómetro es el propio micro:bit. \emph{Si no tienes el aparato físico, el simulador de MakeCode trae una barrita de luz que te deja programar y probar igual;} la medida del cielo real la harán con el dispositivo, aunque sea uno por grupo.}
\milcpaso{2}{milcTurquesa}{\textbf{Pilar 2 --- Un número crudo no dice nada: hay que calibrar.} \texttt{lightLevel} te da \emph{42}. ¿Eso es oscuro o claro? \textbf{Depende de con qué lo compares.} Calibrar es \textbf{fijar umbrales} midiendo condiciones \emph{conocidas}: tapa el sensor (eso es tu «oscuro», anota el número), ponlo bajo la luz del salón (tu «claro», anota), y decide las líneas de corte. Por ejemplo: \emph{menos de 30 = cielo oscuro; de 30 a 90 = intermedio; más de 90 = con luz}. \textbf{Esos números son tuyos y se justifican con tus medidas}, igual que la casera ajustaba su pesa contra un peso patrón. Sin calibrar, tu fotómetro es un número suelto; calibrado, es un instrumento.}
\milcpaso{3}{milcTurquesa}{\textbf{Pilar 3 --- La lógica antes que los bloques: pseudocódigo.} Nunca se empieza arrastrando bloques. Primero se escribe la \textbf{lógica en palabras} (pseudocódigo), como aprendiste en el pensamiento computacional: \emph{``Por siempre: leer luz. Si luz < 30, mostrar luna llena (cielo oscuro). Si no, si luz < 90, mostrar media luna. Si no, mostrar sol (mucha luz). Al oprimir A, mostrar el número exacto''}. \textbf{Esa estructura ---leer → comparar con umbrales → actuar--- es la misma de la curandera} que con su mano distinguía \emph{normal / fiebre leve / fiebre fuerte}. Escribir la lógica antes evita perderse después: los bloques solo \textbf{traducen} lo que ya decidiste.}
\milcpaso{4}{milcTurquesa}{\textbf{Pilar 4 --- Una medida no vale hasta que se prueba.} La casera comprobaba su pesa con un peso patrón; tú compruebas tu fotómetro con \textbf{escenarios de prueba}. Antes de creerle, defines \textbf{5 situaciones} con su resultado \emph{esperado} y anotas el \emph{real}: sensor tapado (espero «oscuro»), bajo el bombillo (espero «con luz»), en penumbra (espero «intermedio»), apuntando a la ventana, apuntando a una pantalla encendida. \emph{Si en algún escenario el instrumento no hace lo esperado, no está listo:} ajustas los umbrales y vuelves a probar. \textbf{Probar no es desconfianza: es lo que convierte un programa en un instrumento en el que se puede creer.}}

\begin{drawbox}{milcTurquesa}{Fotómetro micro:bit + umbrales + pseudocódigo + pruebas}
\textbf{Sensor:} \texttt{input.lightLevel()} $\to$ número de 0 (oscuro) a 255 (mucha luz). Sin cablear. \\[1mm]
\textbf{Calibrar:} medir condiciones conocidas y fijar umbrales (p. ej. <30 oscuro · 30--90 intermedio · >90 con luz). \\[1mm]
\textbf{Pseudocódigo:} leer $\to$ comparar con umbrales $\to$ mostrar. La lógica antes que los bloques. \\[1mm]
\textbf{Prueba:} 5 escenarios con esperado/real; si falla, ajustar umbrales y repetir.
\end{drawbox}

\begin{softbox}{milcMostaza}{milcGris}{Errores comunes y su corrección}
\textbf{Errores que arruinan un fotómetro (y cómo evitarlos).} (1) \emph{Usar umbrales inventados} --- si no mediste tu «oscuro» y tu «claro», tus cortes no significan nada; \textbf{calibra con condiciones reales}. (2) \emph{Arrastrar bloques sin pseudocódigo} --- te pierdes en el código; \textbf{escribe la lógica en palabras primero}. (3) \emph{Confiar sin probar} --- un programa que «corre» no es un instrumento que «mide»; \textbf{pásalo por los 5 escenarios}. (4) \emph{Medir con una pantalla o el propio brillo del micro:bit cerca} --- contaminas la lectura; mide con la matriz \textbf{apuntando limpio} al cielo. (5) \emph{Creer que el número es «la verdad absoluta»} --- el micro:bit compara, no mide en unidades astronómicas; sirve para \textbf{contrastar} lugares y noches. (6) \emph{No anotar la hora y el lugar} --- sin eso, la lectura no se puede comparar, igual que en tu bitácora.
\end{softbox}

\begin{infoband}{milcTurquesa}
\guiaDiagrama{assets/astronomia-3/calibracion-microbit.tex}{Calibrar el fotómetro: el sensor solo da un número (0 a 255); los umbrales lo traducen a significado, y una medida de 42 deja de ser «me pareció oscuro» para volverse un dato comprobable.}

\textbf{Lo que va al cuaderno (Actividad 2 · EXPLICA).} Título: «Actividad 2 --- Mi fotómetro por dentro». \textbf{Formato:} 3 bloques. Bloque A: la barra de calibración 0--255 con mis umbrales marcados y por qué los puse ahí. Bloque B: el pseudocódigo de mi programa (leer → comparar → mostrar). Bloque C: mis 5 escenarios de prueba con esperado/real. \textbf{Extensión:} 1 hoja. \textbf{Sabes que terminaste cuando} podrías darle tu pseudocódigo a otro y armaría el mismo instrumento.
\end{infoband}

\puente{Ya tienes la lógica en la cabeza. Es hora de bajarla a pseudocódigo, armarla en MakeCode y \textbf{ponerla a prueba} ---nadie da por buena una medida sin probarla.}

\milcestacion{milcMagenta}{Estación 3 de 3 · Hacer}

\begin{softbox}{milcMagenta}{milcGris}{Cómo vas a construir tu producto}
\textbf{Actividad 3 · CREA + EVALÚA --- Construye, calibra y contrasta tu fotómetro} (75 min · individual o parejas). Hora de construir. Vas a armar el fotómetro en MakeCode, calibrarlo con tus propias medidas, probarlo en 5 escenarios y ---lo más importante--- \textbf{contrastarlo con tu ojo} del módulo 1.
\end{softbox}

\milcpaso{1}{milcMagenta}{\textbf{Paso 1 --- Calibra primero (15 min).} Antes de programar, \textbf{mide condiciones conocidas} y anótalas. Con un micro:bit (o el simulador): lee \texttt{lightLevel} con el sensor \emph{tapado} (tu «oscuro»: anota el número), bajo la \emph{luz del salón} (tu «claro»: anota), y en \emph{penumbra} (media luz: anota). Con esos tres números, \textbf{decide tus dos umbrales} y escríbelos con su justificación: \emph{``pongo el umbral 1 en \_\_ porque tapado me dio \_\_''}. Sin este paso, todo lo demás es adivinar.}
\milcpaso{2}{milcMagenta}{\textbf{Paso 2 --- Escribe el pseudocódigo (10 min).} En el cuaderno, la lógica en palabras: \emph{``Por siempre: leer luz. Si luz < umbral 1 → mostrar luna (oscuro). Si no, si luz < umbral 2 → mostrar media luna (intermedio). Si no → mostrar sol (con luz). Al oprimir A → mostrar el número''}. Léelo en voz alta: ¿tiene sentido de principio a fin? Corrige aquí, que es barato; corregir en bloques es caro.}
\milcpaso{3}{milcMagenta}{\textbf{Paso 3 --- Ármalo en MakeCode (20 min).} Ahora traduce tu pseudocódigo a bloques en \emph{makecode.microbit.org}. Necesitas: un bucle \textbf{«por siempre»}, la variable \textbf{luz} = \texttt{lightLevel}, una estructura \textbf{si / si no si / si no} con tus umbrales, e íconos distintos para cada rango. Añade: al \textbf{oprimir el botón A}, que \textbf{muestre el número} exacto. Prueba en el simulador moviendo la barrita de luz mientras lo armas.}
\milcpaso{4}{milcMagenta}{\textbf{Paso 4 --- Pruébalo en 5 escenarios (15 min).} Nadie da por buena una medida sin probarla. Llena la tabla: para cada escenario escribe el resultado \textbf{esperado} y luego el \textbf{real}. Escenarios: (1) sensor tapado, (2) bajo el bombillo, (3) en penumbra, (4) apuntando a la ventana, (5) apuntando a una pantalla encendida. \textbf{¿Coincidió esperado con real en los 5?} Si no, \textbf{ajusta un umbral y vuelve a probar} ---y anota qué cambiaste. Ese ajuste es la parte más valiosa.}
\milcpaso{5}{milcMagenta}{\textbf{Paso 5 --- Contrasta con tu ojo del módulo 1 (10 min).} Aquí está el corazón del módulo. Recuerda: \textbf{más luz de fondo = más contaminación = menos estrellas visibles}. Así que donde tu ojo vio \emph{más} estrellas (cielo oscuro), el fotómetro debería dar un número \emph{más bajo}. Haz una \textbf{tabla de contraste}: lugar/noche · estrellas que conté (ojo) · lectura del fotómetro (instrumento) · ¿concuerdan? Escribe 2 frases: dónde el ojo y el instrumento \textbf{coinciden} y dónde \textbf{no}, y una hipótesis de por qué.}
\milcpaso{6}{milcMagenta}{\textbf{Paso 6 --- Sé honesto con tu instrumento (5 min).} Ningún instrumento es perfecto. Escribe \textbf{una cosa en la que le crees} a tu fotómetro (p. ej. distinguir tu patio de una esquina con poste) y \textbf{una en la que no} (p. ej. medir estrellas débiles, para lo que el sensor es muy flojo). \emph{Reconocer el límite de tu propia pesa es lo que te separa de quien cree que su número es la verdad absoluta.}}

\begin{drawbox}{milcMagenta}{Antes de dar por bueno tu fotómetro}
\checkbox Calibré con condiciones conocidas (tapado, penumbra, luz) y anoté los números. \\[1mm]
\checkbox Mis dos umbrales están escritos con su justificación. \\[1mm]
\checkbox Escribí el pseudocódigo antes de armar los bloques. \\[1mm]
\checkbox El programa muestra un ícono por rango y el número al oprimir A. \\[1mm]
\checkbox Probé los 5 escenarios (esperado/real) y ajusté lo que falló. \\[1mm]
\checkbox Hice la tabla de contraste instrumento vs ojo, con 2 frases de análisis.
\end{drawbox}

\begin{softbox}{milcMagenta}{milcGris}{Plantilla de trabajo}
\textbf{Plantilla del fotómetro.} \textit{Calibración:} tapado = \_\_ · penumbra = \_\_ · luz del salón = \_\_. \\[1mm]
\textit{Umbrales:} umbral 1 = \_\_ (porque\ldots) · umbral 2 = \_\_ (porque\ldots). \\[1mm]
\textit{Pseudocódigo:} por siempre → leer luz → si <u1 (oscuro) / si <u2 (intermedio) / si no (con luz) → botón A muestra el número. \\[1mm]
\textit{Pruebas (esperado/real):} tapado \_\_ · bombillo \_\_ · penumbra \_\_ · ventana \_\_ · pantalla \_\_. \\[1mm]
\textit{Contraste:} lugar · estrellas (ojo) · lectura (instrumento) · ¿concuerdan?
\end{softbox}

\begin{infoband}{milcMagenta}


\textbf{Lo que va al cuaderno (Actividad 3 · CREA + EVALÚA).} Título: «Actividad 3 --- Mi fotómetro calibrado». \textbf{Formato:} calibración + umbrales justificados + pseudocódigo + tabla de 5 pruebas + tabla de contraste instrumento vs ojo + 2 frases de análisis + 1 límite reconocido. \textbf{Extensión:} 1-2 hojas. \textbf{Sabes que terminaste cuando} tu fotómetro le daría a otra persona la misma lectura que a ti en el mismo lugar.
\end{infoband}

\puente{El producto no es «un programa que corre»: es un \textbf{instrumento calibrado} que da lecturas repetibles, y una \textbf{tabla que lo contrasta con tu ojo}. Eso ya es medir de verdad.}

\milcestacion{milcMostaza}{Producto de la sesión}
\textbf{Fotómetro micro:bit calibrado + tabla de contraste instrumento vs ojo}.
\begin{softbox}{milcMostaza}{milcGris}{Criterios de calidad}
(1) Los umbrales están calibrados con medidas reales, no inventados. (2) Hay pseudocódigo escrito antes de los bloques, y el programa lo cumple. (3) El programa muestra un ícono por rango y el número exacto al oprimir A. (4) Los 5 escenarios de prueba están documentados (esperado/real) con al menos un ajuste. (5) La tabla de contraste instrumento vs ojo tiene 2 frases de análisis y 1 límite reconocido.
\end{softbox}

\puente{Antes de cerrar, sé honesto con tu instrumento: ¿en qué le crees y en qué no? Un buen investigador conoce los límites de su propia pesa.}

\milcestacion{milcVino}{Evaluación liberadora · cinco dimensiones}

\begin{softbox}{milcVino}{milcGris}{Sentido de la evaluación}
Evaluar no es solo revisar una nota. Es reconocer qué aprendí, qué cambió en mí, qué emociones manejé, cómo me relaciono con la ciudad y la comunidad, y qué le dejaré a quien viene detrás.
\end{softbox}

\textbf{Mi reflexión liberadora en cinco dimensiones.} Completa cada frase iniciada:

\small
\begin{tabularx}{\textwidth}{>{\bfseries\RaggedRight\arraybackslash}p{3.4cm}Y>{\RaggedRight\arraybackslash}p{4.6cm}}
\rowcolor{milcVino}\color{white}Dimensión & \color{white}\bfseries Mi respuesta (frame starter) & \color{white}\bfseries Algo que descubrí\\
1. Personal & Lo que más cambió en mí fue \linead{6.5cm} & \linead{4.2cm}\\
2. Emocional & La emoción más fuerte fue \linead{2.5cm} y la regulé \linead{3cm} & \linead{4.2cm}\\
3. Ciudadana & En democracia esto importa porque \linead{6cm} & \linead{4.2cm}\\
4. Local (Cartago) & En mi barrio sería útil para \linead{6.5cm} & \linead{4.2cm}\\
5. Intergeneracional & A mi abuela/abuelo le diría \linead{6.5cm} & \linead{4.2cm}\\
\end{tabularx}
\normalsize

\begin{infoband}{milcVino}
\textbf{Para la próxima sustentación.} Vuelve a esta tabla en seis meses. Verás que las respuestas cambian — no porque lo aprendido cambie, sino porque tú cambias. La evaluación liberadora es una conversación contigo a través del tiempo.
\end{infoband}


\puente{Cerramos con 3 voces que te ayudan a entender por qué construir un instrumento ---y no solo confiar en el ojo--- es un acto de rigor y de justicia.}

\milcestacion{milcGradeDark}{Triángulo de pensamiento · cierre filosófico}

\begin{softbox}{milcVino}{milcGris}{Dussel · lente del nosotros}
\textit{El otro se revela realmente como otro\ldots como el pobre, el oprimido; el que, a la vera del camino, fuera del sistema, muestra su rostro sufriente y sin embargo desafiante: «¡Tengo hambre!, ¡tengo derecho a comer!»}

Dussel dice que el excluido no pide permiso: exige, y con derecho. El cielo de Cartago está «fuera del sistema» de los grandes observatorios, y por eso casi nadie lo mide. Cuando construyes un fotómetro con tus manos y produces una lectura que otro puede comprobar, le das a tu comunidad con qué reclamar: «aquí hay esta luz de más, y lo puedo demostrar». \textbf{Construir el instrumento es praxis: no te quedas contemplando el problema, lo transformas en evidencia.}

\textbf{¿Construí una herramienta que le da voz y evidencia a mi comunidad, o solo un ejercicio para la nota?}
\end{softbox}
\linead{\linewidth}\\[1mm]
\linead{\linewidth}

\begin{softbox}{milcMostaza}{milcGris}{Estoicismo · Séneca · Cartas a Lucilio, 20 (c. 64 d.C.) · cuidado interior}
\textit{La filosofía enseña a obrar, no a hablar; quiere que cada uno viva a la manera que prescribe, que estén en armonía nuestras palabras con nuestras acciones.}

Séneca desconfía del que solo habla. No basta con \emph{entender} que un instrumento mide mejor que el ojo: hay que \textbf{construirlo, calibrarlo y probarlo}. Un fotómetro que solo existe en tu cabeza no mide nada; uno armado, con sus umbrales justificados y sus cinco pruebas, sí. \textbf{La praxis ---hacer con las manos lo que sabes con la cabeza--- es lo que separa al que opina del que investiga.} Tus palabras («este cielo tiene mucha luz») recién valen cuando tus actos (medir, calibrar, comprobar) las respaldan.

\textbf{¿Me quedé en «entendí cómo funciona», o llegué hasta «lo construí, lo probé y lo puedo demostrar»?}
\end{softbox}
\linead{\linewidth}\\[1mm]
\linead{\linewidth}

\begin{softbox}{milcTurquesa}{milcGris}{Floridi · lente de la infoesfera}
\textit{Las tecnologías digitales están re-ontologizando la naturaleza misma de la infoesfera.}

\emph{Re-ontologizar} suena difícil, pero dice algo simple: una tecnología no solo hace tareas, \textbf{cambia la naturaleza misma de aquello que toca}. Tu fotómetro no le \emph{añade} un número al cielo: transforma tu «me pareció oscuro» (una impresión privada, que se pierde) en un dato público (repetible, comparable, comprobable). El instrumento \textbf{re-ontologiza tu observación}: la vuelve parte de un saber compartido. Por eso construir la herramienta cambia no solo lo que mides, sino lo que \emph{cuenta} como conocer.

\textbf{Mi instrumento, ¿convierte mis impresiones en datos que otros pueden usar, o solo me da números para mí?}
\end{softbox}
\linead{\linewidth}\\[1mm]
\linead{\linewidth}

\begin{infoband}{milcNegro}
\textbf{Nota para el docente.} Las tres formulaciones del triángulo son la síntesis del autor de la guía, no citas textuales de los pensadores. Si vas a citarlos en clase, remite a la obra original.
\end{infoband}

\milcestacion{milcVino}{Mi autoevaluación · marca una casilla por fila}

{\small
\setlength{\extrarowheight}{2.2pt}
\begin{tabularx}{\textwidth}{>{\RaggedRight\arraybackslash\bfseries}p{3.0cm}YYY}
\rowcolor{milcVino}\color{white}\bfseries Criterio & \color{white}\bfseries Lo logré & \color{white}\bfseries En proceso & \color{white}\bfseries Necesito apoyo\\[.6mm]
Comprensión del tema & \milcopcion{Explico las ideas con mis palabras.} & \milcopcion{Reconozco las ideas, falta precisión.} & \milcopcion{Necesito repasar lo básico.}\\[.8mm]
\rowcolor{milcGris}Pensamiento computacional & \milcopcion{Usé los pilares al armar mi producto.} & \milcopcion{Usé algunos pilares.} & \milcopcion{Necesito releer la estación 2.}\\[.8mm]
Aplicación y producto & \milcopcion{Mi producto cumple los criterios.} & \milcopcion{Está armado, con ajustes pendientes.} & \milcopcion{Aún está en construcción.}\\[.8mm]
\rowcolor{milcGris}Reflexión 5 dimensiones & \milcopcion{Respondí cada una con honestidad.} & \milcopcion{Respondí algunas con detalle.} & \milcopcion{Mis respuestas son muy generales.}\\[.8mm]
Triángulo de pensamiento & \milcopcion{Conecté las 3 voces con mi trabajo.} & \milcopcion{Conecté al menos una.} & \milcopcion{No las relaciono con mi proceso.}\\[.8mm]
\end{tabularx}}

\vspace{3mm}
\textbf{Hoy entendí que:}\\[1mm]
\linead{\linewidth}\\[1mm]
\linead{\linewidth}

\textbf{Mi compromiso para la próxima sesión es:}\\[1mm]
\linead{\linewidth}\\[1mm]
\linead{\linewidth}

\begin{softbox}{milcMostaza}{milcGris}{Cita de despedida · Marco Aurelio}
\textit{``El obstáculo en el camino se vuelve el camino. Lo que estorba al camino se vuelve camino.''}\\[1mm]
Cada error de hoy, bien observado, se convierte en pieza del camino que pisarás mañana.
\end{softbox}

\small
\begin{tabularx}{\textwidth}{>{\bfseries}p{3.6cm}Y}
\rowcolor{milcGradeDark!12}Nombre completo: & \linead{10cm}\\
Fecha: & \linead{4cm} \quad Curso/grupo: \linead{4cm}\\
Mi firma: & \linead{9cm}\\
\end{tabularx}
\normalsize

\begin{infoband}{milcGradeDark}
\textbf{Para la próxima sesión --- mide el cielo con tu instrumento.} Esta es la tarea que cierra el contraste ojo/instrumento. \textbf{Cuándo:} las mismas \textbf{3 noches} (o 3 momentos) y a la \textbf{misma hora}, para poder comparar. \textbf{Qué haces:} en cada punto, apunta el fotómetro \textbf{limpio} hacia el cielo (sin pantallas ni el propio brillo cerca), oprime A y \textbf{anota el número}, junto con \textbf{lugar y hora}. Mide al menos en \textbf{2 lugares distintos} (tu patio y una esquina más oscura, o con y sin poste cerca). \textbf{Luego, completa tu tabla de contraste} con tu magnitud límite a ojo del módulo 1: donde contaste más estrellas, ¿el fotómetro dio número más bajo? \textbf{Trae:} el instrumento, la tabla de lecturas y \textbf{una conclusión en una frase} sobre cuánta luz de más tiene tu barrio. \emph{En el módulo 4 vas a comunicar todo esto: la observación a ojo, los datos, el instrumento y lo que hallaste ---y a reconocer sus límites y su ética.}
\end{infoband}

\end{document}
